\section{Forest and Tree Models}
\label{sec:extensionplus}


In this section, we extend our tableau calculus to ensure that extracted models are \emph{data trees}.
Besides their theoretical interest, date trees are relevant as they are the underlying structure of XML documents in practice.
In this extension of the calculus, we take advantage of the expressive power of nominals and the possibility to ``jump'' to evaluation points to characterize \emph{forests}, i.e., disjoint unions of trees, and trees as a particular case.
We introduce rules and clash conditions to determine whether nodes have at most one predecessor, and to avoid loops via accessibility relations.  
We then extend our soundness and completeness results to deal with the
class of forests.
Finally, we show how to  get completeness with respect to the class of trees by
checking a single constraint on open branches.

To work with forests and trees, we extend our tableaux rules with a new type of accessibility constraint called {\em transitivity constraint}.
Transitivity constraints are expressions of the form $i{:}\tup{+} j$.
They are introduced into the calculus to represent reachability navigating accessibility relations.
We introduce transitivity constraints and the new rules of the calculus in~\Cref{rules:reachability}.  
In this figure, the rule (+intro) ensures that the relation described by ${+}$ contains the relation described by ${\pathh}$, whereas the rule (trans) indicates that the relation described by ${+}$ is a transitive relation.
Finally, the rule (+copy) takes care of dealing with urfathers, copying transitivity constraints where they need to be copied.
 
\begin{figure}[t]% REACHABILITY RULES
    \begin{center}
        \small 
        % \resizebox*{\textwidth}{!}
        {\def\arraystretch{1.2}
        \begin{tabular}{c}
            \toprule
            {\sc Reachability Rules}
            \tabularnewline
            \midrule
            {
            \begin{tabular}{ccc}
                {\infer[(+\mathrm{intro})]
                { i{:}\tup{+} j}
                { i{:}\tup{\pathh} j}
                }
                &
                {\infer[(\mathrm{trans})]
                    {i{:} \tup{+} k}
                    {\deduce[]
                    { j{:}\tup{+} k}
                    {i{:}\tup{+} j} 
                    }
                }
                &
                {\infer[(+\mathrm{copy})]
                    {i{:}\tup{+} k}
                    {\deduce[]
                    {\phantom{|} j{:}k} 
                    {\phantom{|} i{:}\tup{\pathh} j}}
                }
            \end{tabular}
            }
            \tabularnewline
            \bottomrule                        
        \end{tabular}}
    \end{center}
    \caption{Tableau Rules for Forests and Trees.}
    \label{rules:reachability}
\end{figure}


\begin{rem}
    % It is important to point out that transitivity constraints are not part of the language of $\Hxpath$.
    % They are auxiliary annotations for the tableaux.
     It is important to point out that even though transitivity constraints can be mistaken as characterizing the transitive closure of $\bigcup_{\pathh \in \rels} R_{\pathh}$, this is not the case. The relation associated to $\tup{+}$ in the extracted model may not be the smallest transitive relation containing $\bigcup_{\pathh \in \rels} R_{\pathh}$.
    In any case, transitivity constraints suffice for handling forests (and trees as a particular case).
    
    Furthermore, we can always assume that the symbol $+$ is not part of the set $\rels$, and then does not occur in the node expression at the root of the tableau. Thus, the obtained tableau is still internalized, but over the extended signature $\rels'=\rels\cup\{+\}$.
\end{rem}

% For a sake of simplicity, in what follows we consider the fragment of $\Hxpath$ with a single relational symbol and its ``reachability'' counterpart, i.e., $\rels=\{\pathh,+\}$.

In addition to the rules in~\Cref{rules:reachability}, we introduce new clash conditions for forests and trees.

\begin{defi}\label{def:clashf}
    A branch $\Theta$ is said to contain a \emph{clash} if it satisfies any of the conditions listed in \Cref{def:clash}, or if any of the following situations arise:
    \begin{enumerate}
        \setcounter{enumi}{2}
        \item $i{:}\tup{+}i \in \Theta$, 
        \item $\{j{:}\tup{+} i, k{:}\tup{+} i\} \subseteq \Theta$, and $\{j{:}k,j{:}\tup{+}k\}\cap\Theta = \emptyset$. % $j\neq k$. 
    \end{enumerate}
\end{defi}

Notice how condition (3) in \Cref{def:clashf} is an indication that no node can be reached by itself; whereas condition (4) in \Cref{def:clashf} tells us that a node cannot have more than one predecessor. 
More precisely, condition (4) tells us that if a nominal is reachable from two others, then, they either indicate the same node, or one must be reachable from the other.


The following example demonstrates how the introduction of new rules and clash conditions restricts the range of models that satisfy a given node expression.

\begin{exa}
    \Cref{fig:btableau} depicts a tableau for $\tup{@_0 \pathh {0?} =_{\cmp} p?}$.
    In brief, $\tup{@_0\pathh 0? =_{\cmp} p?}$ forces a self loop at the node named by $0$.
    Applying only the rules in \Cref{rules:basic,rules:nominals,rules:internalization,rules:path} from \Cref{sec:tableaux} yields the first 11 nodes of the tableau.
    Notice that if we restrict our attention only to those 11 nodes (i.e., if we truncate the tableau at the node 11), we obtain a branch that is open and saturated according to \Cref{def:saturated} (modulo the application of reflexivity and symmetry rules, which are omitted).
    This indicates that $\tup{@_0\pathh 0? =_{\cmp} p?}$ is satisfiable in the class of all models.
    The nodes 12 and 13 show the effect of the new rules, and the way in which they allow us to talk about reachability.
    Notice also that node 13 introduces a clash condition as per \Cref{def:clashf}; i.e., a node is reachable from itself.
    This clash condition tells that $\tup{@_0 \pathh {0?} =_{\cmp} p?}$ is unsatisfiable in the class of tree models; and thus also unsatisfiable in the class of forests.
    \begin{figure}[!h]
        \begin{center}
        \begin{prooftree}{}%
        [
            1{:}\tup{@_0 \pathh {0?} =_{\cmp} p?},
            just={root}
            [
                \tup{@_1@_0 \pathh {0?} =_{\cmp} @_1p?},
                just={(int$_1$):!u}
                [
                    \tup{@_0 \pathh {0?} =_{\cmp} @_1p?},
                    just={($@$):!u}
                    [
                        0{:}\tup{\pathh}2,
                        just={(child):!u}
                        [
                            \tup{@_2{0?} =_{\cmp} @_1p?},
                            just={(child):!uu}
                            [
                                2{:}0,
                                just={(test):!u}
                                [
                                    \tup{@_2 =_{\cmp} @_1p?},
                                    just={(test):!uu}
                                    [
                                        \tup{@_1p? =_{\cmp} @_2},
                                        just={(com$_1$):!u}
                                        [
                                            1{:}p,
                                            just={(test):!u}
                                            [
                                                \tup{@_1 =_{\cmp} @_2},
                                                just={(test):!uu}
                                                [
                                                    \tup{@_2 =_{\cmp} @_1},
                                                    just={(com$_1$):!u}
                                                    [%
                                                        % 0{:}0,
                                                        % just={(ref):!uuuuuuuuuuu}
                                                        % [
                                                        %     1{:}1,
                                                        %     just={(ref):!uuuuuuuuu}
                                                        %     [
                                                        %         2{:}2,
                                                        %         just={(ref):!uuuuuuuuuu}
                                                        %         [
                                                        %             \tup{@_0 =_{\cmp} @_0},
                                                        %             just={(dRef):!uuuuuuuuuuuuuu}
                                                        %             [
                                                        %                 \tup{@_1 =_{\cmp} @_1},
                                                        %                 just={(dRef):!uuuuuuuuuuuuuuu}
                                                        %                 [
                                                        %                     \tup{@_2 =_{\cmp} @_2},
                                                        %                     just={(dRef):!uuuuuuuuuuuuu}
                                                        %                     [
                                                        %                         0{:}2,
                                                        %                         just={(sym):!uuuuuuuuuuuu}
                                                        %                         [
                                                                                    0{:}\tup{+}2,
                                                                                    just={($+$intro):!uuuuuuuu}
                                                                                    [
                                                                                        0{:}\tup{+}0,
                                                                                        just={($+$copy):!u,!uuuuuuu},
                                                                                        close={:!u}
                                                                                    ]
                                                        %                         ]
                                                        %                     ]
                                                        %                 ]
                                                        %             ]
                                                        %         ]
                                                        %     ]
                                                        % ]
                                                    ]
                                                ]
                                            ]
                                        ]
                                    ]
                                ]
                            ]
                        ]
                    ]
                ]
            ]
        ]
        \end{prooftree}
        \end{center}
       \caption{A closed tableau for $\tup{@_0 \pathh {0?} =_{\cmp} p?}$}\label{fig:btableau}
    \end{figure}
\end{exa}

We are now in a position to prove that the new rules and clash conditions indeed force forest-like structures.

\begin{prop}
\label{prop:forest}
Let $\Theta$ be an open saturated branch in the extended calculus, and let $\model^{\Theta}$ be its extracted model.  
Then, $\model^{\Theta}$ is a forest.	
\end{prop}

\begin{proof}
    The proof is concluded if $F = (\bigcup_{\pathh\in\rels} R_{\pathh}^{\Theta})^{+}$ is a forest.
    To this end, we show that:
    \begin{enumerate}[(a)]
    \item it is not the case that $(n,n) \in \mathbin{F}$, and
    % \item for all $j\in\cdom$, there are $\pathh,\astep[b]\in\rels$ and $i\in\cdom$ such that $j\cra i$ and $k R^{\Theta}_{\astep[b]} i$, then $j=k$.
    \item if $(n,m)\in \mathbin{F}$ and $(n',m) \in \mathbin{F}$, then $n = n'$ or $(n,n') \in \mathbin{F}$.
    \end{enumerate} 

    The proof of (a) is by contradiction.
    Suppose that $(n,n)\in \mathbin{F}$.
    Then, exist $\astep[a]_1,\ldots,\astep[a]_{\ell}$ such that
        $(n,i)\in R_{\astep[a]_1}, (i,j)\in R_{\astep[a]_2} \ldots (k,n)\in R_{\astep[a]_{\ell}}$.
    From the definition of $\cra$, we have that ${\{n{:}\tup{\pathh_1} i, i{:}\tup{\pathh_2}j,\ldots,k{:}\tup{\pathh_{\ell}} n\}\subseteq\Theta}$.
    Since $\Theta$ is closed under (+intro), we also have that $\{n{:}\tup{+} i, i{:}\tup{+}j,\ldots,k{:}\tup{+} n\}\subseteq\Theta$.
    In turn, closure under (trans), gives us ${n{:}\tup{+} n}\in \Theta$.
    This yields a clash using condition (3) from~\Cref{def:clashf}, contradicting the assumption that $\Theta$ is open.

    The proof of (b) is also by contradiction.
    Suppose $(n,m)\in \mathbin{F}$, $(n',m) \in \mathbin{F}$, $n \neq n'$, and $(n,n')\notin\mathrel{F}$.
    From $(n,m)\in \mathbin{F}$ and $(n',m)\in \mathbin{F}$, we know that $\{n{:}\tup{+}m,n'{:}\tup{+}m\} \subseteq \Theta$.
    In addition, from $(n,n')\notin \mathrel{F}$, we know $n{:}\tup{+}n' \notin \Theta$.
    Using the definition of urfather, we know that $j{:}k \notin \Theta$.
    When put together, these facts yield a clash using condition (4) from~\Cref{def:clashf}, contradicting the assumption that $\Theta$ is open.
\end{proof}

To further ensure that the resulting model is a tree we only need 
to verify that every point in the extracted model is reachable from 
the root. 

\begin{defi}
    \label{def:connected}
    Let $\Theta$ be a branch of a tableau for $i{:}\varphi$ in the extended calculus. We say that $\Theta$
    is {\em connected} iff for all $j\in\Theta$, if $i{:}j \notin \Theta$, then, $i{:}\tup{+}j \in \Theta$.
\end{defi}

\begin{prop}
    \label{prop:connectedtree}
    Let $\Theta$ be an open, saturated, and connected branch. Then, $\model^{\Theta}$ is a tree. 
\end{prop}

\begin{proof}
    The proof is concluded if $T = (\bigcup_{\pathh\in\rels} R_{\pathh}^{\Theta})^{+}$ is a tree.
    We already know that $T$ is a forest.
    Thus, we only need to prove that for all $j\in\cdom$, $(i,j)\in \mathbin{T}$.
    This is straightforward since $\Theta$ is connected.
\end{proof}

We conclude this section by showing that the satisfiability problem for forests and data trees remains decidable and stays within \pspace.

\begin{thm}
    \label{th:terminationtrees}
   % Any tableau in the extended calculus is finite.  
   The satisfiability problem for forests and data trees in $\Hxpath$ is decidable and \pspace-complete.
\end{thm}

\begin{proof}
    Notice that for forests and trees, we only add three new rules which are applied over
    a finite set of node expressions (by~\Cref{prop:termination}). Since we apply rules only once, this yields again a finite tableau. Moreover, the new rules are clearly of type 1, thus applying them in~\Cref{alg:sat} yields similar bounds as in~\Cref{th:satpspace}. Checking the clash conditions can also be done using polynomial space. Therefore, the extended algorithm runs in polynomial space.
\end{proof}

